\subsection*{G01 -- Fossile Energiewirtschaft \normalfont\small (21 Firmen, Archetyp A1)}
\noindent Der Fußabdruck steckt im Produkt, nicht im Betrieb. Trotzdem ist die betriebliche Seite gut messbar und trennscharf: Methanschlupf und Abfackelung unterscheiden Förderer um Faktoren, während ihre Produktemissionen fast identisch sind.

\noindent\textbf{Materielle Kennzahlen}
\begin{itemize}[nosep,leftmargin=1.2em]
  \item Methanintensität der Förderung
  \item Abfackelungsrate
  \item Scope-1-Intensität je Barrel Öläquivalent
  \item Reserven-Exposure (stranded assets)
  \item Investitionsanteil in emissionsarme Energie
\end{itemize}
\noindent\textbf{Bezugsgröße:} Barrel Öläquivalent bzw. verarbeitete Menge, nicht Umsatz (Umsatz schwankt mit dem Ölpreis und verzerrt Zeitreihen)\\
\textbf{Datenlage:} hoch: EPA GHGRP und Climate TRACE erfassen Anlagen dieser Gruppe praktisch vollständig\\
\textbf{Untergruppen (Ebene C):} Förderung, Midstream, Raffination, Ausrüster\\
\textbf{Mitglieder:} {\raggedright\small\ttfamily APA, BKR, COP, CVX, DVN, EOG, EQT, EXE, FANG, HAL, KMI, MPC, OKE, OXY, PSX, SLB, TPL, TRGP, VLO, WMB, XOM\par}

\medskip
\subsection*{G02 -- Versorger \& Umweltinfrastruktur \normalfont\small (33 Firmen, Archetyp A1)}
\noindent Regulierte Netzmonopole mit eigener Verbrennung. Der Erzeugungsmix erklärt fast die gesamte Streuung, deshalb ist er die Kennzahl und nicht der Umsatz. Abfall- und Wasserwirtschaft steht hier, weil Deponiemethan derselben Logik folgt: eigene Anlage, behördlich gemeldet, physisch messbar.

\noindent\textbf{Materielle Kennzahlen}
\begin{itemize}[nosep,leftmargin=1.2em]
  \item CO2-Intensität je erzeugter MWh
  \item Kohleanteil an der Erzeugung und Stilllegungsplan
  \item Deponiemethan und Gasverwertungsquote
  \item Anteil erneuerbarer Kapazität am Zubau
  \item Netzverluste bzw. Wasserverluste
\end{itemize}
\noindent\textbf{Bezugsgröße:} MWh Erzeugung, Tonne Abfall, Kubikmeter Wasser\\
\textbf{Datenlage:} hoch: EPA GHGRP deckt Kraftwerke und Deponien über der 25.000-Tonnen-Schwelle ab\\
\textbf{Untergruppen (Ebene C):} Stromerzeuger, Mehrspartenversorger, Wasser, Abfall\\
\textbf{Mitglieder:} {\raggedright\small\ttfamily AEE, AEP, AES, ATO, AWK, CEG, CMS, CNP, D, DTE, DUK, ED, EIX, ES, ETR, EVRG, EXC, FE, LNT, NEE, NI, NRG, PCG, PEG, PNW, PPL, RSG, SO, SRE, VST, WEC, WM, XEL\par}

\medskip
\subsection*{G03 -- Rohstoff- \& Prozessindustrie \normalfont\small (26 Firmen, Archetyp A1)}
\noindent Chemie, Zement, Stahl, Papier: ein Teil der Emissionen ist chemisch unvermeidbar (Kalzinierung, Reduktion) und nicht durch Grünstrom zu beseitigen. Diese Gruppe darf man nur gegen sich selbst messen, sonst bestraft das Modell Physik statt Verhalten.

\noindent\textbf{Materielle Kennzahlen}
\begin{itemize}[nosep,leftmargin=1.2em]
  \item Scope-1-Intensität je Tonne Produkt
  \item Energieintensität der Prozesse
  \item Wasserentnahme in Wasserstressgebieten
  \item Gefährlicher Abfall und Luftschadstoffe
  \item Anteil Sekundär- bzw. Rezyklatmaterial
\end{itemize}
\noindent\textbf{Bezugsgröße:} Tonne Produkt (Umsatz je Tonne schwankt mit Rohstoffpreisen)\\
\textbf{Datenlage:} hoch: klassische GHGRP-Melder\\
\textbf{Untergruppen (Ebene C):} Chemie, Baustoffe, Metalle \& Bergbau, Papier \& Verpackung\\
\textbf{Mitglieder:} {\raggedright\small\ttfamily ALB, AMCR, APD, AVY, BALL, CF, CRH, CTVA, DD, DOW, ECL, FCX, IFF, IP, LIN, LYB, MLM, MOS, NEM, NUE, PKG, PPG, SHW, STLD, SW, VMC\par}

\medskip
\subsection*{G04 -- Transport \& Logistik \normalfont\small (16 Firmen, Archetyp A2)}
\noindent Flottenbetreiber: Fluggesellschaften, Reedereien, Bahnen, Speditionen. Der Treibstoffverbrauch ist der Fußabdruck, und er ist je Leistungseinheit direkt vergleichbar. Kreuzfahrtreedereien stehen hier und nicht bei den Hotels: ein Schiff ist ein Verbrennungsmotor mit Betten.

\noindent\textbf{Materielle Kennzahlen}
\begin{itemize}[nosep,leftmargin=1.2em]
  \item Treibstoffintensität je Tonnen- bzw. Passagierkilometer
  \item Flottenalter und -effizienz
  \item Anteil alternativer Treibstoffe (SAF, LNG, Elektrifizierung)
  \item Luftschadstoffe und Schwefelgehalt (Schifffahrt)
  \item Arbeitssicherheit
\end{itemize}
\noindent\textbf{Bezugsgröße:} Tonnenkilometer, verfügbare Sitzkilometer, Bettennächte\\
\textbf{Datenlage:} mittel: Climate TRACE modelliert Luft- und Seeverkehr, EPA erfasst ortsfeste Anlagen nur teilweise\\
\textbf{Untergruppen (Ebene C):} Luftfahrt, Schifffahrt, Schiene, Straße\\
\textbf{Mitglieder:} {\raggedright\small\ttfamily CCL, CHRW, CSX, DAL, EXPD, FDX, FDXF, JBHT, LUV, NCLH, NSC, ODFL, RCL, UAL, UNP, UPS\par}

\medskip
\subsection*{G05 -- Fahrzeug-, Luftfahrt- \& Wehrtechnik \normalfont\small (17 Firmen, Archetyp A3)}
\noindent Langlebige Investitionsgüter, deren Emissionen beim Kunden entstehen. Ein Autohersteller mit sauberer Fabrik und Verbrennerflotte ist nicht nachhaltig; die Werksemission zu ranken wäre hier das klassische Fehlmaß.

\noindent\textbf{Materielle Kennzahlen}
\begin{itemize}[nosep,leftmargin=1.2em]
  \item Flottenemissionswert der verkauften Produkte (g CO2/km, kg/Flugstunde)
  \item Anteil emissionsfreier Antriebe am Absatz
  \item Materialintensität und Recyclingfähigkeit
  \item Lücke zwischen Reduktionsversprechen und Ist-Entwicklung
  \item Arbeitssicherheit in der Fertigung
\end{itemize}
\noindent\textbf{Bezugsgröße:} verkaufte Einheit bzw. Produktflotte, nicht Umsatz\\
\textbf{Datenlage:} keine kostenlose Quelle für die Nutzungsphase; Flottenwerte nur aus Unternehmensberichten und Zulassungsstatistiken\\
\textbf{Untergruppen (Ebene C):} Automobil, Zivile Luftfahrt, Verteidigung, Zulieferer\\
\textbf{Mitglieder:} {\raggedright\small\ttfamily APTV, AXON, BA, F, GD, GE, GM, HII, HONA, HWM, LHX, LMT, NOC, RTX, TDG, TSLA, TXT\par}

\medskip
\subsection*{G06 -- Investitionsgüter \& Bauwirtschaft \normalfont\small (48 Firmen, Archetyp A3)}
\noindent Maschinen-, Elektro- und Bauindustrie. Diskrete Fertigung mit moderatem eigenem Ausstoß, relevanter Vorkette (Stahl, Kupfer) und Wirkung über die Effizienz des verkauften Geräts. Bauträger und Bauunternehmen stehen hier, weil ihr Hebel ebenfalls im Gebauten liegt, nicht im eigenen Betrieb.

\noindent\textbf{Materielle Kennzahlen}
\begin{itemize}[nosep,leftmargin=1.2em]
  \item Scope-1- und Scope-2-Intensität der Fertigung
  \item Energieeffizienz bzw. Emissionsvermeidung des verkauften Produkts
  \item Anteil Umsatz aus Effizienz- und Elektrifizierungslösungen
  \item Eingebundener Kohlenstoff (Beton, Stahl) bei Bau und Bauträgern
  \item Arbeitssicherheit
\end{itemize}
\noindent\textbf{Bezugsgröße:} Umsatz (bereinigt) je Segment; bei Bauträgern je Wohneinheit\\
\textbf{Datenlage:} mittel: einzelne Werke in GHGRP, Produktkennzahlen nur aus Berichten\\
\textbf{Untergruppen (Ebene C):} Maschinenbau, Elektrotechnik, Bauprodukte, Bau \& Bauträger\\
\textbf{Mitglieder:} {\raggedright\small\ttfamily ALLE, AME, AOS, BLDR, CARR, CAT, CMI, DE, DHI, DOV, EME, EMR, ETN, FAST, FERG, FIX, FTV, GEV, GNRC, GWW, HON, HUBB, IEX, IR, ITW, J, JCI, LEN, LII, MAS, MMM, NDSN, NVR, OTIS, PCAR, PH, PHM, PNR, PWR, ROK, SNA, SWK, TT, URI, VLTO, VRT, WAB, XYL\par}

\medskip
\subsection*{G07 -- Halbleiter \& Elektronikfertigung \normalfont\small (45 Firmen, Archetyp A4)}
\noindent Die Gruppe mit der größten internen Spreizung und deshalb die wichtigste Tier-Unterscheidung: ein Fab-Betreiber trägt Prozessgase (hohes GWP), gigantischen Strom- und Wasserbedarf selbst; ein Fabless-Designer und ein Markenhersteller ohne eigene Fertigung haben denselben Fußabdruck in der Vorkette. Wer beide auf einer Skala misst, belohnt Auslagerung.

\noindent\textbf{Materielle Kennzahlen}
\begin{itemize}[nosep,leftmargin=1.2em]
  \item Scope-1- und Scope-2-Intensität (nur Fab-Betreiber)
  \item Prozessgase mit hohem Treibhauspotenzial (PFC, SF6, NF3)
  \item Wasserentnahme und Wiederverwendungsquote
  \item Auftragsfertigungs-Programm und Lieferanten-Energiepolitik (Fabless)
  \item Produktenergieeffizienz und Elektroschrott-Rücknahme
\end{itemize}
\noindent\textbf{Bezugsgröße:} Wafer-Startfläche bzw. Umsatz; Tier getrennt normalisieren\\
\textbf{Datenlage:} Fab-Betreiber mittel (US-Werke in GHGRP), Fabless und Markenhersteller niedrig -- Climate TRACE erfasst diesen Sektor nicht\\
\textbf{Untergruppen (Ebene C):} Fab-Betreiber, Fabless, Ausrüstung \& Material, Geräte \& Markenhardware\\
\textbf{Mitglieder:} {\raggedright\small\ttfamily AAPL, ADI, AMAT, AMD, ANET, APH, AVGO, CIEN, COHR, CSCO, DELL, FFIV, FLEX, FSLR, GLW, HPE, HPQ, INTC, JBL, KEYS, KLAC, LITE, LRCX, MCHP, MPWR, MRVL, MSI, MU, NTAP, NVDA, NXPI, ON, Q, QCOM, ROP, SMCI, SNDK, STX, SWKS, TDY, TEL, TER, TXN, WDC, ZBRA\par}

\medskip
\subsection*{G08 -- Nahrungsmittel, Getränke \& Tabak \normalfont\small (19 Firmen, Archetyp A4)}
\noindent Agrarvorkette dominiert: Landnutzung, Methan aus Tierhaltung, Düngemittel, Wasser. Wasserverbrauch ist hier existenziell und für ein Softwarehaus bedeutungslos -- das Lehrbuchbeispiel für Materialität.

\noindent\textbf{Materielle Kennzahlen}
\begin{itemize}[nosep,leftmargin=1.2em]
  \item Wasserentnahme, gewichtet nach Wasserstress am Standort
  \item Landnutzungs- und Entwaldungspolitik der Beschaffung
  \item Verpackungsmaterial und Rezyklatanteil
  \item Lebensmittelverluste
  \item Scope-1-/2-Intensität der Verarbeitung
\end{itemize}
\noindent\textbf{Bezugsgröße:} Tonne bzw. Hektoliter Produkt; Umsatz nur als Ersatzgröße\\
\textbf{Datenlage:} mittel: Verarbeitungswerke teils in GHGRP, Agrarvorkette kostenlos nicht belastbar\\
\textbf{Untergruppen (Ebene C):} Verarbeitung, Getränke, Agrarhandel, Tabak\\
\textbf{Mitglieder:} {\raggedright\small\ttfamily ADM, BF.B, BG, GIS, HRL, HSY, KDP, KHC, KO, MDLZ, MKC, MNST, MO, PEP, PM, SJM, STZ, TAP, TSN\par}

\medskip
\subsection*{G09 -- Markenkonsumgüter (Non-Food) \normalfont\small (14 Firmen, Archetyp A4)}
\noindent Haushalt, Körperpflege, Bekleidung, Schuhe: eigene Fertigung klein, Wirkung in Auftragsfertigung, Chemie, Textilfärberei und Verpackung. Sozialkennzahlen der Lieferkette sind hier nicht Beiwerk, sondern das Hauptrisiko.

\noindent\textbf{Materielle Kennzahlen}
\begin{itemize}[nosep,leftmargin=1.2em]
  \item Lieferanten-Audits und Arbeitsbedingungen in der Vorkette
  \item Wasser- und Chemikalieneinsatz der Auftragsfertigung
  \item Materialherkunft und Rezyklatanteil
  \item Verpackung und Rücknahme
  \item Scope-1-/2-Intensität der eigenen Standorte
\end{itemize}
\noindent\textbf{Bezugsgröße:} Umsatz; Einheiten nur segmentweise vergleichbar\\
\textbf{Datenlage:} niedrig: keine Anlagendaten, Auswertung aus Berichten (HuggingFace-Korpus) und Auditlisten\\
\textbf{Untergruppen (Ebene C):} Haushalt \& Pflege, Bekleidung \& Schuhe, Freizeitgüter\\
\textbf{Mitglieder:} {\raggedright\small\ttfamily CHD, CL, CLX, DECK, EL, GRMN, HAS, KMB, KVUE, LULU, NKE, PG, RL, TPR\par}

\medskip
\subsection*{G10 -- Handel, Distribution \& Gastronomie \normalfont\small (31 Firmen, Archetyp A4)}
\noindent Fläche, Kühlung, Sortiment. Der größte eigene Hebel sind Kältemittelverluste -- eine Kennzahl, die in Gesamtscores meist untergeht, obwohl sie bei Lebensmittelhändlern zweistellige Prozentanteile des Scope 1 ausmacht. Gastronomie steht hier, weil Beschaffung und Kühlung dieselbe Struktur haben.

\noindent\textbf{Materielle Kennzahlen}
\begin{itemize}[nosep,leftmargin=1.2em]
  \item Kältemittelverluste (CO2-Aequivalent je Filiale bzw. Fläche)
  \item Energieintensität je Quadratmeter Verkaufs- und Lagerfläche
  \item Emissionen der eigenen Auslieferungsflotte
  \item Lebensmittel- und Verpackungsabfall
  \item Beschaffungspolitik für kritische Rohstoffe
\end{itemize}
\noindent\textbf{Bezugsgröße:} Quadratmeter Fläche bzw. Filialzahl, zusätzlich Umsatz\\
\textbf{Datenlage:} niedrig bis mittel: Kältemittel nur aus Berichten, Verteilzentren teils in GHGRP\\
\textbf{Untergruppen (Ebene C):} Lebensmittelhandel, Fachhandel, Distribution, Gastronomie\\
\textbf{Mitglieder:} {\raggedright\small\ttfamily AMZN, AZO, BBY, CAH, CASY, CMG, COR, COST, CVNA, DG, DLTR, DPZ, DRI, GPC, HD, HSIC, KR, LOW, MCD, MCK, ORLY, ROST, SBUX, SYY, TGT, TJX, TSCO, ULTA, WMT, WSM, YUM\par}

\medskip
\subsection*{G11 -- Pharma, Biotech \& Medizintechnik \normalfont\small (42 Firmen, Archetyp A4)}
\noindent Wirkstoffproduktion ist Feinchemie mit Lösemittel- und Sonderabfallproblem; Medizintechnik ist diskrete Fertigung mit Einwegprodukten und Sterilisationschemie (Ethylenoxid). Gemeinsam ist die regulierte Produktsicherheit als Governance-Achse.

\noindent\textbf{Materielle Kennzahlen}
\begin{itemize}[nosep,leftmargin=1.2em]
  \item Scope-1-/2-Intensität der Produktion
  \item Gefährlicher Abfall und Lösemittelbilanz
  \item Wirkstoffeinträge in Gewässer
  \item Sterilisationsemissionen und Einwegmaterialanteil
  \item Produktsicherheit und Rückrufe
\end{itemize}
\noindent\textbf{Bezugsgröße:} Umsatz; bei Wirkstoffherstellern zusätzlich Tonne Wirkstoff\\
\textbf{Datenlage:} mittel: Chemieanlagen teils in GHGRP, Entsorgungsdaten in US-Behördenregistern\\
\textbf{Untergruppen (Ebene C):} Pharma, Biotech, Labortechnik, Medizingeräte\\
\textbf{Mitglieder:} {\raggedright\small\ttfamily A, ABBV, ABT, ALGN, AMGN, BAX, BDX, BIIB, BMY, BSX, COO, CRL, DHR, DXCM, EW, GEHC, GILD, IDXX, INCY, ISRG, JNJ, LLY, MDT, MRK, MRNA, MTD, PFE, PODD, REGN, RMD, RVTY, SOLV, STE, SYK, TECH, TMO, VRTX, VTRS, WAT, WST, ZBH, ZTS\par}

\medskip
\subsection*{G12 -- Gesundheitsversorgung, Kostenträger \& klinische Dienste \normalfont\small (12 Firmen, Archetyp A5)}
\noindent Kliniken, Labore, Dialyse, Auftragsforschung, Krankenversicherer. Kein Produkt, keine Fabrik: der Fußabdruck ist Gebäudeenergie plus Narkosegase, und bei den Kostenträgern liegt er faktisch im Kapitalanlageportfolio. Die Gruppe braucht deshalb eine eigene Kennzahlenwahl.

\noindent\textbf{Materielle Kennzahlen}
\begin{itemize}[nosep,leftmargin=1.2em]
  \item Energieintensität je Quadratmeter Klinik- und Laborfläche
  \item Narkose- und Kältegase
  \item Medizinischer Sonderabfall
  \item Kapitalanlage-Exposure der Versicherer (analog PCAF)
  \item Patientensicherheit und Datenschutz
\end{itemize}
\noindent\textbf{Bezugsgröße:} Quadratmeter bzw. Belegungstag; bei Versicherern Prämienvolumen\\
\textbf{Datenlage:} niedrig: praktisch keine kostenlose Anlagendatenquelle\\
\textbf{Untergruppen (Ebene C):} Kostenträger, Leistungserbringer, Labore \& Forschung\\
\textbf{Mitglieder:} {\raggedright\small\ttfamily CI, CNC, CVS, DGX, DVA, ELV, HCA, HUM, IQV, LH, UHS, UNH\par}

\medskip
\subsection*{G13 -- Software \& digitale Plattformen \normalfont\small (34 Firmen, Archetyp A5)}
\noindent Anlagenarm: kein eigenes Rechenzentrum, kein Werk. Der Fußabdruck ist eingekaufte Cloud-Leistung, Büroenergie und Dienstreisen -- also fast vollständig Scope 3. Genau hier sind Gesamtscores am irreführendsten, weil eine niedrige absolute Emission wie Leistung aussieht, obwohl sie Geschäftsmodell ist.

\noindent\textbf{Materielle Kennzahlen}
\begin{itemize}[nosep,leftmargin=1.2em]
  \item Emissionen eingekaufter Cloud-Leistung (Scope 3 Kat. 1)
  \item Büroenergie je Mitarbeitendem und Dienstreisen
  \item Datenschutz und Datensicherheit
  \item Humankapital: Fluktuation, Löhne, Diversität
  \item Governance: Aktionärsrechte, Stimmrechtsstruktur, Vergütung
\end{itemize}
\noindent\textbf{Bezugsgröße:} Mitarbeitende (die einzige belastbare physische Bezugsgröße in dieser Gruppe)\\
\textbf{Datenlage:} niedrig: weder EPA noch Climate TRACE erfassen diese Gruppe; Governance- und Sozialkennzahlen aus SEC-Einreichungen\\
\textbf{Untergruppen (Ebene C):} Unternehmenssoftware, Marktplätze \& Plattformen, IT-Dienste\\
\textbf{Mitglieder:} {\raggedright\small\ttfamily ABNB, ACN, ADBE, ADSK, BKNG, CDNS, CDW, CRM, CRWD, CSGP, CTSH, DASH, DDOG, EBAY, EXPE, FICO, FTNT, GEN, IBM, INTU, IT, LDOS, NOW, PANW, PLTR, PTC, RDDT, SNPS, TRMB, TTWO, TYL, UBER, VEEV, WDAY\par}

\medskip
\subsection*{G14 -- Rechenzentren, Netze \& Funktürme \normalfont\small (19 Firmen, Archetyp A5)}
\noindent Betreiber physischer digitaler Infrastruktur: Hyperscaler, Telekommunikation, Rechenzentrums- und Funkturm-REITs. Sie sind Großverbraucher von Strom und Kühlwasser und damit untereinander hervorragend vergleichbar -- aber nicht mit anlagenarmer Software. Dass GICS Microsoft und einen 300-Personen-Softwareanbieter in dieselbe Sub-Industry legt, ist der teuerste Vergleichsfehler im Technologiesektor.

\noindent\textbf{Materielle Kennzahlen}
\begin{itemize}[nosep,leftmargin=1.2em]
  \item Energieverbrauch und Effizienz (PUE) der Rechenzentren
  \item Anteil erneuerbarer Energie und Beschaffungsqualität (stundenscharf vs. bilanziell)
  \item Wasserverbrauch der Kühlung, gewichtet nach Wasserstress
  \item Notstrom-Scope-1 (Diesel) und Kältemittel
  \item Elektroschrott und Hardware-Lebensdauer
\end{itemize}
\noindent\textbf{Bezugsgröße:} MWh Verbrauch bzw. installierte Kapazität (MW), nicht Umsatz\\
\textbf{Datenlage:} mittel: Verbrauch aus Unternehmensberichten belastbar, Emissionsfaktoren aus Netzdaten ableitbar\\
\textbf{Untergruppen (Ebene C):} Hyperscaler, Telekommunikation, Rechenzentrums-Vermieter\\
\textbf{Mitglieder:} {\raggedright\small\ttfamily AKAM, AMT, CCI, CHTR, CMCSA, DLR, ECHO, EQIX, GDDY, GOOG, GOOGL, META, MSFT, ORCL, SBAC, T, TMUS, VRSN, VZ\par}

\medskip
\subsection*{G15 -- Transaktions- \& Informationsdienstleister \normalfont\small (21 Firmen, Archetyp A5)}
\noindent Zahlungsverkehr, Börsen, Ratings, Datenanbieter. Sie sehen wie Finanzunternehmen aus, tragen aber kein Kreditportfolio: es gibt keine finanzierten Emissionen zu messen. Ihr Fußabdruck ist Rechenzentrumsstrom, ihr Hauptrisiko ist Governance -- bei Ratingagenturen und Indexanbietern zusätzlich der Interessenkonflikt, den sie selbst bewerten.

\noindent\textbf{Materielle Kennzahlen}
\begin{itemize}[nosep,leftmargin=1.2em]
  \item Energieintensität der Transaktionsinfrastruktur
  \item Anteil erneuerbarer Energie
  \item Datensicherheit und Systemverfügbarkeit
  \item Governance und Interessenkonflikte
  \item Humankapital
\end{itemize}
\noindent\textbf{Bezugsgröße:} Transaktionsvolumen bzw. Mitarbeitende\\
\textbf{Datenlage:} niedrig: keine Anlagendaten\\
\textbf{Untergruppen (Ebene C):} Zahlungsverkehr, Marktinfrastruktur, Daten \& Ratings\\
\textbf{Mitglieder:} {\raggedright\small\ttfamily BR, CBOE, CME, COIN, CPAY, EFX, FDS, FIS, FISV, GPN, ICE, JKHY, MA, MCO, MSCI, NDAQ, PYPL, SPGI, V, VRSK, XYZ\par}

\medskip
\subsection*{G16 -- Banken, Versicherer \& Vermögensverwalter \normalfont\small (52 Firmen, Archetyp A6)}
\noindent Die einzige Gruppe, in der eigene Emissionen strukturell bedeutungslos sind: laut CDP liegen über 99 Prozent im Portfolio. Eine Bank nach Büroenergie zu ranken ist die sauberste Ausblendung, die ein ESG-Score anbieten kann -- sie misst Klimaanlagen statt Kreditbücher.

\noindent\textbf{Materielle Kennzahlen}
\begin{itemize}[nosep,leftmargin=1.2em]
  \item Finanzierte Emissionen nach PCAF (Kredit- und Anlagebuch)
  \item Exposure gegenüber fossiler Förderung und Kohleverstromung
  \item Versicherungstechnisches Exposure (Underwriting fossiler Projekte)
  \item Stimmrechtsausübung und Engagement-Politik
  \item Governance: Vergütung, Aufsichtsstruktur, Fehlverhaltensfälle
\end{itemize}
\noindent\textbf{Bezugsgröße:} Bilanzsumme, verwaltetes Vermögen, Prämienvolumen\\
\textbf{Datenlage:} niedrig: PCAF-Werte nur aus freiwilligen Berichten; Stimmrechtsdaten aus SEC N-PX öffentlich\\
\textbf{Untergruppen (Ebene C):} Banken, Versicherer, Vermögensverwalter, Konsumentenkredit\\
\textbf{Mitglieder:} {\raggedright\small\ttfamily ACGL, AFL, AIG, AIZ, ALL, AMP, APO, ARES, AXP, BAC, BEN, BLK, BNY, BRK.B, BX, C, CB, CFG, CINF, COF, EG, FITB, GL, GS, HBAN, HIG, HOOD, IBKR, IVZ, JPM, KEY, KKR, L, MET, MS, MTB, NTRS, PFG, PGR, PNC, PRU, RF, RJF, SCHW, STT, SYF, TFC, TROW, TRV, USB, WFC, WRB\par}

\medskip
\subsection*{G17 -- Immobilien \& Gebäudebetrieb \normalfont\small (28 Firmen, Archetyp A5)}
\noindent Bestandshalter und Betreiber von Fläche: REITs, Hotelbetreiber, Casinoresorts. Vergleichbar über Energie je Quadratmeter -- eine Kennzahl, die es in dieser Gruppe tatsächlich flächendeckend gibt (GRESB, Zertifizierungen), anders als in den meisten anderen.

\noindent\textbf{Materielle Kennzahlen}
\begin{itemize}[nosep,leftmargin=1.2em]
  \item Energieintensität je Quadratmeter (kWh/m2)
  \item Anteil zertifizierter Fläche und Sanierungsquote
  \item Eingebundener Kohlenstoff bei Neubau
  \item Wasserverbrauch und Klimarisiko-Exposure der Standorte
  \item Kältemittel
\end{itemize}
\noindent\textbf{Bezugsgröße:} Quadratmeter bzw. Zimmer, nicht Umsatz\\
\textbf{Datenlage:} mittel: Flächen- und Verbrauchsdaten in REIT-Berichten vergleichsweise standardisiert\\
\textbf{Untergruppen (Ebene C):} Wohnen, Gewerbe, Beherbergung, Spezial\\
\textbf{Mitglieder:} {\raggedright\small\ttfamily ARE, BXP, CPT, DOC, ESS, EXR, FRT, HLT, HST, INVH, IRM, KIM, LVS, MAA, MAR, MGM, O, PLD, PSA, REG, SPG, UDR, VICI, VMRK, VTR, WELL, WY, WYNN\par}

\medskip
\subsection*{G18 -- Medien \& unternehmensnahe Dienste \normalfont\small (25 Firmen, Archetyp A5)}
\noindent Inhalte, Werbung, Personal, Versicherungsmakler, Facility-Dienste. Anlagenarm wie G13, aber ohne Cloud-Kostentreiber; der Fußabdruck sind Produktionen, Veranstaltungen, Fuhrparks und Büros. Eigene Gruppe, weil ihre relevanten Achsen sozial und governancebezogen sind -- nicht energetisch.

\noindent\textbf{Materielle Kennzahlen}
\begin{itemize}[nosep,leftmargin=1.2em]
  \item Büroenergie und Fuhrpark bzw. Produktionsemissionen
  \item Dienstreisen und Veranstaltungslogistik
  \item Humankapital: Arbeitsbedingungen, Fluktuation, Löhne
  \item Governance: Stimmrechtsstruktur, Unabhängigkeit des Aufsichtsrats
  \item Integrität der Inhalte bzw. Beratungsunabhängigkeit
\end{itemize}
\noindent\textbf{Bezugsgröße:} Mitarbeitende\\
\textbf{Datenlage:} niedrig: keine Anlagendaten, Auswertung aus Berichten\\
\textbf{Untergruppen (Ebene C):} Medien \& Inhalte, Werbung, Personal \& Beratung, Facility\\
\textbf{Mitglieder:} {\raggedright\small\ttfamily ADP, AJG, AON, APP, BRO, CBRE, CPRT, CTAS, DIS, ERIE, FOX, FOXA, LYV, MRSH, NFLX, NWS, NWSA, OMC, PAYX, PSKY, ROL, TKO, TTD, WBD, WTW\par}
